\section{Introduction}
\label{sec:introduction}

A sigmoid transformer can exactly realize belief propagation
computations. This is not a loose analogy or a high-level functional
similarity: under the constructive mapping proved in the companion
paper~\cite{coppola2026transformers}, the transformer forward pass
implements the BP update equations directly. The present paper is a
tutorial guide to what that result means, how it connects to
transformer components, and what it implies for understanding trained
models.

The formal anchor is the companion theorem. This paper does not
rederive it. Instead, it builds the conceptual picture around it:
opening each transformer component, identifying its role in the
message-passing account, and connecting the formal result to a
converging body of empirical work in mechanistic interpretability,
in-context learning, and graphical model theory.

\subsection*{Scope of the Claims}

The paper operates at four distinct levels:

\begin{enumerate}
\item \textbf{Exact constructive results.} The companion theorem gives
an exact realization of BP in sigmoid transformers under a constructive
weight mapping~\cite{coppola2026transformers}.

\item \textbf{Empirical observations.} Mechanistic interpretability
results from trained transformers exhibit structures consistent with
the BP interpretation.

\item \textbf{Mechanistic interpretation.} Transformer components can
be understood through a routing-and-inference decomposition aligned
with BP primitives.

\item \textbf{Speculative extensions.} The discussions of grounding,
hallucination, and hybrid continuous-discrete inference are
interpretive extensions rather than established theorems.
\end{enumerate}

These claims have different evidentiary status and are marked
accordingly throughout the paper.

\subsection*{The Five Correspondences}

The tutorial is organized around five structural correspondences
between transformer components and the operations of belief
propagation:

\begin{enumerate}
\item \textbf{The residual stream is the belief state.} Each token's
$D_{\mathrm{model}}$-dimensional vector stores the current accumulated
state of inference at that node. Residual connections preserve and
update beliefs across layers rather than replacing them outright.

\item \textbf{Attention performs evidence gathering.} Attention heads
route information from neighboring tokens into the residual stream
before the update step executes. Multiple required evidence sources
become jointly available in a shared workspace.

\item \textbf{The FFN performs belief update.} Under the constructive
BP mapping, the feed-forward network computes
$\updateBelief(m_0, m_1) = \sigma(\logit(m_0) + \logit(m_1))$ from
the gathered evidence. This is the Bayesian update rule for
independent binary evidence sources --- exact with sigmoid activation,
approximate with ReLU or SwiGLU.

\item \textbf{Layers correspond to rounds of message passing.} One
layer is one gather-then-update cycle. Stacking layers corresponds
naturally to iterative rounds of inference.

\item \textbf{The scaling dimensions describe inference capacity.} The
architectural dimensions --- $W$, $L$, $H$, $D_{ff}$,
$D_{\mathrm{model}}$ --- determine the scale of the routing and
inference operations available to the network.
\end{enumerate}

\subsection*{The Running Example}

Throughout the paper we will repeatedly use a simple example involving
three propositions:
\[
\texttt{like(jack,jill)}, \quad \texttt{like(jill,jack)}, \quad
\texttt{date(jack,jill)}.
\]

The example is intentionally small. Its purpose is not realism but
clarity. We use it repeatedly to illustrate how attention gathers
evidence, how FFNs combine beliefs, how residual streams store
intermediate state, how layers propagate information over multiple
rounds, and what grounding means in a formally declared concept space.
By keeping the example fixed across sections, the paper can focus on
computational structure rather than constantly introducing new
notation.

\subsection*{The Underlying Algebra}

All five correspondences rest on the same mathematical spine: the
log-odds algebra of independent binary evidence, developed by Turing
and Good~\cite{good1950probability}, formalized by Pearl as the
sum-product algorithm~\cite{pearl1988probabilistic}, and realized
exactly by the constructive sigmoid mapping. Within this mapping,
sigmoid is the exact activation required to recover the log-odds
update equations. Part~I develops this foundation before turning to
the architectural correspondences.

\subsection*{Grounded and Ungrounded Inference}

The distinction between grounded and ungrounded transformers is not
primarily architectural. In both cases, the network performs inference
over an implicit factor graph. The difference is whether the graph has
a declared ontology, finite concept space, and verifier that permits
formal correctness guarantees.

For grounded transformers operating on tree-structured knowledge
bases, the constructive BP results provide exact posterior guarantees.
Standard language models do not provide guarantees of this kind
because their internal representations are not tied to a formally
declared ontology. Part~IV develops this distinction and its
implications for hallucination and reliability.

\subsection*{Convergent Evidence}

Five independent research programs have converged on the same general
picture from different starting points: the transformer forward pass
implements a form of distributed probabilistic inference. Mechanistic
interpretability studies identify routing heads and structured FFN
updates. Belief-state geometry work finds posterior-like organization
in residual streams. In-context learning work models transformers as
Bayesian learners. Graph neural network research established the
message-passing connection. Constructive weight-learning experiments
show that optimization reliably discovers BP-like solutions. Part~V
surveys this convergence.

\subsection*{How to Read This Paper}

Part~I establishes the mathematical foundation. Readers already
familiar with log-odds algebra and belief propagation can skim
portions of it. Part~II develops the core architectural
correspondences component by component. Part~III examines
superposition, the three softmaxes, and the relationship between
trained weights and implicit factor graphs. Part~IV addresses
grounding and hallucination. Part~V surveys convergent evidence and
related work.

The formal constructive results are developed
in~\cite{coppola2026transformers}. This paper is a tutorial guide to
the conceptual landscape those results open up.